Dim ultraviolet light of wavelength \laO\ and intensity \IaO\ falls on a clean zinc plate (work function \WaO). Electrons leave the plate as soon as the light is switched on.
\begin{abcliste}
\abc In the wave model, the energy of the light arrives continuously and is spread evenly over the surface. Assume that one electron collects the energy falling on a square of side \daO\ (about the size of an atom). What power does it receive?
\abc How long would the electron have to wait, in this model, before it had collected the work function? Compare with the observation.
\abc Explain the observation with photons: calculate the energy of one photon in eV and the number of photons per second that hit \AaO\ of the plate.
\end{abcliste}
